%%%%%%%%%%%%%%%%%%%%%%%
%%% DOCUMENT SETTINGS
%%%%%%%%%%%%%%%%%%%%%%%
\documentclass[11pt]{exam}
\printanswers % Uncomment to show answers
%%%%%%%%%%%%%%
%%% PACKAGES
%%%%%%%%%%%%%%
\usepackage{amsmath,amsfonts,amssymb,amsthm} % math fonts and symbols
\usepackage{bbm} % Math blackboard font
\usepackage{enumitem} % For reference enumerations lists
\usepackage{textcomp} % some symbols (like copyright)
\usepackage[top = 2cm, bottom = 2cm, right = 1.2cm, left = 1.2cm, paper = a4paper, headheight = 6cm]{geometry} % Margins setup
\usepackage{xcolor} % Colors
\usepackage{graphicx}
%%%%%%%%%%%%%%%%%%%%%%%%%%
%%% MACROS AND COMMANDS
%%%%%%%%%%%%%%%%%%%%%%%%%%
\newcommand{\N}{\mathbb{N}} % natural number set
\newcommand{\R}{\mathbb{R}} % real number set
\makeatletter
\renewcommand{\@makefnmark}{\textsuperscript{\arabic{footnote}}} % Ensure numeric superscripts
\renewcommand{\thefootnote}{\arabic{footnote}} % Ensure numeric labels
\makeatother
\setcounter{footnote}{1}
\newcommand{\BAR}{%
    \hspace{-\arraycolsep}%
    \strut\vrule % the `\vrule` is as high and deep as a strut
    \hspace{-\arraycolsep}%
}
%%%%%%%%%%%%%%%%%%%%%%%%%%%%%%%%%
%%% HEADER AND FOOT 
%%%%%%%%%%%%%%%%%%%%%%%%%%%%%%%%%
% Defining information
\newcommand{\degree}{Bachelor in Philosophy, Politics and Economics}
\newcommand{\shortdeg}{PPE}
\newcommand{\exam}{1st Midterm - Retake}
\newcommand{\subject}{Mathematics}
\newcommand{\theyear}{2025/26}
\newcommand{\ccauthor}{A. G. Azze - CUNEF Universidad}
\newcommand{\thedate}{Oct 10, 2025}
\newcommand{\university}{CUNEF Universidad}
% Setting header and footer
\pagestyle{headandfoot}
\header{\degree \\ \subject\ - \theyear}{}{\thedate \\ \ccauthor}
\footer{\university}{}{\thepage}
%%%%%%%%%%%%%%%%%%%%%%%%%
%%% DOCUMENT CONTENT
%%%%%%%%%%%%%%%%%%%%%%%%%
\begin{document}
%--- Instructions ---%
\begin{center}

    {\Huge \textbf{\exam{}}}
    \vspace{0.7cm}

    \begin{flushright}
        \textbf{Time Limit:} 50 minutes
    \end{flushright}
    \vspace{-0.2cm}
    
    \fbox{\parbox{\textwidth}{\centering 
        \begin{enumerate}[leftmargin = 0.5cm, label = $\bullet$, rightmargin = 0.2cm]
    
            \item \textbf{Every non-trivial calculation and claim in the exam must be properly justified.}

            \item The use of calculators is not necessary and therefore not allowed.

            \item Digital devices such as mobile phones, tablets, and laptops, as well as internet access, are forbidden.
            
        \end{enumerate}
    }}
\end{center}
%--- Instructions ---%

%--- Questions ---%

\begin{questions}

% Question
\question[3]
Consider the demand and supply functions
\[
p_D(q)=20 - q^2 + 2q,
\qquad
p_S(q)=\frac{q^2 - 14q + 70}{5}.
\]

\begin{enumerate}[leftmargin=0.5cm,label=(\alph*),rightmargin=0.2cm]
  \item Find the supply–demand equilibrium.
  \item Let $M(q)=p_D(q)-p_S(q)$ be the per–unit margin. What value of $q$ maximizes $M(q)$?
\end{enumerate}

%% Solution
\begin{solution}[0cm]
\begin{enumerate}[leftmargin=0.5cm,label=(\alph*),rightmargin=0.2cm]

  \item Solve $p_D(q)=p_S(q)$:
  \[
  M(q)=p_D(q)-p_S(q)
  =-\frac{6}{5}\bigl(q^2-4q-5\bigr)
  =-\frac{6}{5}(q-5)(q+1).
  \]
  Hence $(q-5)(q+1)=0\Rightarrow q\in\{-1,5\}$. The feasible equilibrium is
  \[
  q^*=5,\qquad p^*=p_D(5)=20-25+10=5.
  \]

  \item Since $M(q)=-\dfrac{6}{5}(q^2-4q-5)$, then
\[
M'(q)=-\frac{6}{5}\,(2q-4)=-\frac{12}{5}(q-2),\qquad
M''(q)=-\frac{12}{5}<0.
\]
The only critical point ($M'(q)=0$) is $q=2$. Since $M'(q)=-\dfrac{12}{5}(q-2)$, the only critical point is $q=2$. Moreover,
$M''(q)=-\dfrac{12}{5}<0$ for all $q$, so $M$ is concave everywhere and $q=2$
gives the global maximum. The maximum value is
\[
M(2)=-\frac{6}{5}\,(4-8-5)=\frac{54}{5}.
\]
\end{enumerate}
\end{solution}





% Question
\question[2]
    The biodiversity of a forest decays at $2\%$ every year. How long will it take for the biodiversity to be halved?    
%% Solution
\begin{solution}[0cm]
Let $B(t)$ represent the biodiversity at year $t$. Then, $B(t)=B_0(0.98)^t$. We want to find $t$ such that $B(t)=\tfrac{1}{2}B_0$:
\[
(0.98)^t=\tfrac12
\quad\Longrightarrow\quad
t=\frac{\ln(1/2)}{\ln(0.98)}\approx \frac{-0.6931}{-0.0202}\approx 34.3\ \text{years}.
\]
\end{solution}

% Question
\question[2]
    Define the functions $f_1(q) = 50 - \frac{q^2}{q^3+1}$ and $f_2(q) = 5(10 + e^{-q})$. 
    \begin{enumerate}[leftmargin=0.5cm,label=(\alph*),rightmargin=0.2cm]
        \item What is the domain of $f_1(q) = 50 - \frac{q^2}{q^3+1}$? What is the domain of $f_2(q) = 5(10 + e^{-q})$
        \item Find the limits $\lim_{q\rightarrow \infty} f_1(q)$ and $\lim_{q\rightarrow \infty} f_2(q)$.
        \item Define another function $h(q)$ such that, for all $q > 0$,
        \[
        f_1(q) \;\le\; h(q) \;\le\; f_2(q).
        \]
        Compute the limit $\lim_{q\rightarrow \infty} h(q)$ (or at least find upper and lower bounds for the limit)
    \end{enumerate}
    
%% Solution
\begin{solution}[0cm]
\begin{enumerate}[leftmargin=0.5cm,label=(\alph*),rightmargin=0.2cm]

  \item 
  \[
  f_1(q)=50-\frac{q^2}{q^3+1}\quad\Rightarrow\quad q^3+1\neq 0\ \Longrightarrow\ q\neq -1.
  \]
  Hence \(\mathrm{Dom}(f_1)=\mathbb{R}\setminus\{-1\}\).
  Since \(f_2(q)=5\big(10+e^{-q}\big)\) is defined for every \(q\in\mathbb{R}\), we have
  \(\mathrm{Dom}(f_2)=\mathbb{R}\).

  \item 
  \begin{align*}
      \lim_{q\to\infty} f_1(q)
      &= \lim_{q\to\infty}\!\left(50-\frac{q^2}{q^3+1}\right)
      = 50-\lim_{q\to\infty}\frac{q^2}{q^3(1+1/q^3)}
      = 50-\lim_{q\to\infty}\frac{1}{q(1+1/q^3)}
      = 50. \\
      \lim_{q\to\infty} f_2(q) 
      &= \lim_{q\to\infty} 5\big(10+e^{-q}\big)=5(10+0)=50.
    \end{align*}

  \item Since for all $q>0$ we have $f_1(q)\le h(q)\le f_2(q)$, taking limits as $q\to\infty$ and using part (b):
\[
\lim_{q\to\infty} f_1(q)=50,\qquad \lim_{q\to\infty} f_2(q)=50.
\]
Hence,
\[
\lim_{q\to\infty} f_1(q)\ \le\ \lim_{q\to\infty} h(q)\ \le\ \lim_{q\to\infty} f_2(q),
\]
so by the Squeeze Theorem, $\displaystyle \lim_{q\to\infty} h(q)=50$.

\end{enumerate}
\end{solution}


% Question
\question[3]
Let $q(p)$ be the number of units a company sells when the price-per-unit is $p$, let $p(t)$ be the price per unit at time $t$. Define $Q(t) = q(p(t))$.
\begin{enumerate}[leftmargin = 0.5cm, label = (\alph*), rightmargin = 0.2cm]
    \item Provide an interpretation of $Q(t)$.
    \item Provide an approximation of $Q(1.5)$, given that $q(20) = 50$, $q'(20) = 10$, $p(1) = 20$, and $p'(1) = 1$.
    \item Let $R(t) = Q(t)p(t)$ be the revenue of the company at time $t$. \\
    Assume that $R(t)$ reaches a minimum at $t^*$. 
    \begin{enumerate}[leftmargin = 0.5cm, label = (\roman*), rightmargin = 0.2cm]
        
        \item Prove that $\frac{Q'(t^*)}{Q(t^*)} = -\frac{p'(t^*)}{p(t^*)}$.

        \item What is the sign of $R'(t)$ for values of $t$ near $t^*$?  
        
    \end{enumerate}
\end{enumerate}

% Solution
\begin{solution}[0cm]
\begin{enumerate}[leftmargin=0.5cm,label=(\alph*),rightmargin=0.2cm]
  \item $Q(t)=q(p(t))$ is the quantity sold at time $t$, induced by the price $p(t)$ set at time $t$.

  \item Linear approximation at $t=1$:
  \[
  Q'(1) = q'(p(1))\,p'(1)=q'(20)\cdot 1=10,\quad 
  Q(1) = q(20) = 50, \quad
  Q(1.5)\approx Q(1)+Q'(1)\cdot 0.5=55.
  \]

  \item $R(t)$ has a minimum at $t^*$.
    \begin{enumerate}[leftmargin = 0.5cm, label = (\roman*), rightmargin = 0.2cm]
      \item With $R(t)=Q(t)p(t)$,
      \[
      R'(t)=Q'(t)p(t)+Q(t)p'(t).
      \]
      At a minimum $t^*$, $R'(t^*)=0$, so
      \[
      Q'(t^*)\,p(t^*)+Q(t^*)\,p'(t^*)=0
      \;\Longrightarrow\;
      \boxed{\;\dfrac{Q'(t^*)}{Q(t^*)}=-\,\dfrac{p'(t^*)}{p(t^*)}\;}.
      \]
      \item Since $t^*$ is a local minimum, $R'(t)\leq0$ for $t<t^*$ and $R'(t)\geq0$ for $t>t^*$ (i.e., it changes from non-positive to non-negative across $t^*$).
    \end{enumerate}
\end{enumerate}
\end{solution}

\end{questions}	

\newpage
\begin{center}
You may find useful the following derivatives of basic functions and rules of derivatives
\end{center}

\begin{center}
    \textbf{Differentiation rules}
\end{center}
        
        \begin{enumerate}[leftmargin = 0.5cm, label = , rightmargin = 0.2cm]

        \item \textbf{Linearity:} $h(x) = a f(x) + b g(x) \Longrightarrow h'(x) = a f'(x) + b g'(x)$, for $a, b \in\R$

        \item \textbf{Product rule:} $h(x) = f(x)g(x) \Longrightarrow h'(x) = f'(x)g(x) + f(x)g'(x)$

        \item \textbf{Quotient rule:} $h(x) = f(x)/g(x) \text{ and } g(x) \neq 0 \Longrightarrow h'(x) = \frac{f'(x)g(x) - f(x)g'(x)}{g(x)^2}$

        \item \textbf{Chain rule:} $h(x) = g(f(x)) \Longrightarrow h'(x) = g'(f(x))f'(x)$

        \item \textbf{Inverse rule:} $h(x) = f^{-1}(y) \Longrightarrow h'(y) = 1/f'(f^{-1}(y))$
        
        \end{enumerate}

\begin{center}
    \textbf{Derivative of usual functions}
\end{center}

        \begin{enumerate}[leftmargin = 0.5cm, label = , rightmargin = 0.2cm]

        \item \textbf{Constant:} $f(x) = a,\; a \in\R \Longrightarrow f'(x) = 0$
        
        \item \textbf{Power:} $f(x) = x^a \Longrightarrow f'(x) = a x^{a-1}$
        

        \item \textbf{Exponential:} $f(x) = e^x \Longrightarrow f'(x) = e^x$

        \item \textbf{Logarithm:} $f(x) = \ln(x) \Longrightarrow f'(x) = 1/x$

        \item \textbf{Trigonometric:} $f(x) = \sin(x),\; g(x) = \cos(x) \Longrightarrow f'(x) = \cos(x),\; g'(x) = -\sin(x)$
        
        \end{enumerate}

\end{document}
